% ============================================================
%  Chapter 6 — Prime Numbers on the Helix
% ============================================================
\section{Prime Numbers on the Helix}\label{sec:primes}

The conical helix provides a natural geometric embedding of the
congruence classes $n \equiv 1, 5 \pmod{6}$. Since these are exactly the
classes containing all primes greater than~$3$
(Theorem~\ref{thm:primes-mod6}), the primes appear as a distinguished
subset of the integer helix points.

\subsection{Integer Points and the Prime Corollary}

By Corollary~\ref{cor:theta-param}, the helix intersects the plane
$z = n$ at the angle $\theta_n = \frac{\pi}{3}n + \frac{\pi}{4}$
for all $n \equiv 1, 5 \pmod{6}$. Combining this with the classical
fact that all primes $p > 3$ lie in these congruence classes yields
the central geometric corollary:

\begin{theorem}[Primes on the helix]\label{thm:primes-helix}
All primes $p > 3$ lie as discrete points on the conical helix at angles
\begin{equation}\label{eq:prime-angle}
  \boxed{\theta_p = \frac{\pi}{3}\,p + \frac{\pi}{4}.}
\end{equation}
\end{theorem}

\begin{proof}
\textbf{Step~1.} Every prime $p > 3$ satisfies
$p \equiv 1$ or $5 \pmod{6}$ (Theorem~\ref{thm:primes-mod6}).

\textbf{Step~2.} The helix attains integer height $z = n$ precisely at
$\theta_n = \frac{\pi}{3}n + \frac{\pi}{4}$ for
$n \equiv 1, 5 \pmod{6}$ (Corollary~\ref{cor:theta-param}).

\textbf{Conclusion.} Since every prime $p > 3$ satisfies the condition
from Step~1, it lies on the helix at $\theta_p$.
\end{proof}

\begin{example}[First primes on the helix]
\begin{center}
\begin{tabular}{cccc}
\toprule
$p$ & $\theta_p$ (exact) & $\theta_p$ (numerical) & $k(p)$ \\
\midrule
 5 & $23\pi/12$ &  6.021 &  11 \\
 7 & $31\pi/12$ &  8.116 &  20 \\
11 & $47\pi/12$ & 12.305 &  46 \\
13 & $55\pi/12$ & 14.399 &  63 \\
17 & $71\pi/12$ & 18.588 & 105 \\
19 & $79\pi/12$ & 20.682 & 130 \\
23 & $95\pi/12$ & 24.871 & 188 \\
\bottomrule
\end{tabular}
\end{center}
\end{example}

\subsection{Composite Numbers and Complete Characterisation}

The helix does not contain only primes. Since it embeds all indices
$n \equiv 1, 5 \pmod{6}$, composite numbers from these classes also
appear. The first composite on the helix is $n = 25 = 5^2$.

\begin{remark}[Complete characterisation]
The helix contains exactly the set
$\{n \in \mathbb{Z}_{>0} : n \equiv 1, 5 \pmod{6}\}$.
The primes $p > 3$ form a distinguished but non-isolable subset.
The numbers $2$, $3$, and all $n$ with $n \equiv 0, 2, 3, 4 \pmod{6}$
do not lie on the helix.
\end{remark}

\subsection{Angular Spacings and Prime Gaps}

For consecutive primes $p < q$ with $p, q > 3$, the angular gap
$\theta_q - \theta_p = \frac{\pi}{3}(q-p)$ depends only on the
prime gap $g = q - p$. Two particularly noteworthy cases:

\begin{center}
\begin{tabular}{lcc}
\toprule
Prime type & Gap $g$ & Angular spacing \\
\midrule
Twin primes ($q = p+2$) & 2 & $2\pi/3$ \\
Cousin primes ($q = p+4$) & 4 & $4\pi/3$ \\
Sexy primes ($q = p+6$) & 6 & $2\pi$ (full revolution) \\
\bottomrule
\end{tabular}
\end{center}

Twin primes thus occupy adjacent positions on the same $120°$-arc,
while sexy primes are vertically aligned on the helix.

\subsection{Asymptotic Density}

\begin{proposition}[Density of primes on the helix]\label{prop:density}
The density of primes among the integer helix points tends to zero:
\[
  \frac{\#\{p \leq N : p \text{ prime},\; p > 3\}}
       {\#\{n \leq N : n \equiv 1, 5 \pmod{6}\}}
  \;\sim\; \frac{N/\ln N}{N/3}
  \;=\; \frac{3}{\ln N}
  \;\xrightarrow{N\to\infty}\; 0.
\]
\end{proposition}

\subsection{Comparison with Other Prime Spirals}

\begin{remark}[Derived vs.\ empirical]
In contrast to the Ulam spiral and the Sacks spiral, which represent
empirically discovered patterns, the present helix is \emph{derived}
algebraically from the quadratic sequence $k(n)$. The appearance of all
primes $p > 3$ is a \emph{theorem} (Theorem~\ref{thm:primes-helix}),
not an observation. The underlying mechanism --- that the congruence
classes $1, 5 \pmod{6}$ are simultaneously the integrality locus of
$k(n)$ and the residue classes coprime to~$6$ --- is an algebraic
coincidence rooted in the factorisation
$2n^2+3n+1 = (2n+1)(n+1)$.
\end{remark}

\subsection{Open Questions}

\begin{enumerate}
  \item Can the helix structure yield insights into prime gaps or the
    distribution of twin primes?
  \item Does the arc-length parametrisation via $Q(k)$ provide a
    natural metric on the primes?
  \item Which other quadratic sequences lead to analogous helix
    constructions, and what are their geometric invariants?
\end{enumerate}
